\documentclass[conference]{IEEEtran}
\IEEEoverridecommandlockouts

\usepackage{cite}
\usepackage{amsmath,amssymb,amsfonts}
\usepackage{algorithmic}
\usepackage{graphicx}
\usepackage{textcomp}
\usepackage{xcolor}
\usepackage{booktabs}
\usepackage{multirow}
\usepackage{siunitx}
\usepackage[caption=false,font=footnotesize]{subfig}
\usepackage{hyperref}


\newcommand{\acDRL}{deep reinforcement learning (DRL)}
\newcommand{\acHVAC}{heating, ventilation, and air conditioning (HVAC)}
\newcommand{\acRL}{reinforcement learning (RL)}
\newcommand{\acMDP}{Markov decision process (MDP)}
\newcommand{\acOOD}{out-of-distribution (OOD)}
\newcommand{\acCOP}{coefficient of performance (COP)}
\newcommand{\acGHI}{global horizontal irradiance (GHI)}
\newcommand{\acMAE}{mean absolute error (MAE)}
\newcommand{\acRMSE}{root-mean-square error (RMSE)}
\newcommand{\acDR}{domain randomisation (DR)}
\newcommand{\acSB}{Stable-Baselines3 (SB3)}

\graphicspath{{figures/}{../results/plots/}}

\title{DRL-Based HVAC Control for a Single-Zone Smart Building:\\%
A Comparative Study of Agents%
\thanks{Midterm Report submitted for ML709: IoT Smart Systems, Services,
and Applications. Instructor: Prof.~Mohsen Guizani, Mohamed Bin Zayed
University of Artificial Intelligence.}}

\author{%
  \IEEEauthorblockN{%
    Saif Alkindi\IEEEauthorrefmark{1},~%
    Hamdan Alhadhrami\IEEEauthorrefmark{1},~%
    Khalfan Hableel\IEEEauthorrefmark{1}~and~%
    Ali Alharthi\IEEEauthorrefmark{1}%
  }
  \IEEEauthorblockA{%
    \IEEEauthorrefmark{1}Mohamed Bin Zayed University of Artificial Intelligence\\%
    Student IDs: 25011027,~25010875,~25011026,~22010335%
  }%
}

\begin{document}

\maketitle

\begin{abstract}
Buildings consume a large share of global energy, and HVAC is a major load. Efficient HVAC control is hard because weather, prices, and thermal dynamics vary over time. This paper studies deep reinforcement learning for continuous HVAC control in a single-zone smart building. We compare PPO, A2C, SAC, and TD3 in a physics-based environment with nonlinear heat transfer, a Carnot-based coefficient of performance, monthly electricity prices, domain randomization, and balanced seasonal sampling. Agents are trained on Syracuse weather and tested on held-out Albany weather to measure geographic generalization. Over 100 evaluation episodes, TD3 performs best with \SI{0.12}{\degreeCelsius} MAE, \SI{16.9}{\kilo\watt\hour\per day} energy use, and $99.7\%$ of time within $\pm\SI{0.5}{\degreeCelsius}$ of setpoint. A2C and SAC also generalize well, while PPO performs worse. These results show that a simple simulator can support strong performance on unseen weather when training includes seasonal diversity and parameter uncertainty.
\end{abstract}

\begin{IEEEkeywords}
Deep reinforcement learning, HVAC control, smart buildings, energy efficiency,
PPO, SAC, TD3, A2C, domain randomisation
\end{IEEEkeywords}

\section{Introduction}
\label{sec:introduction}
Buildings remain a major source of energy use, and heating, ventilation, and air-conditioning (HVAC) systems account for a large share of that demand. In practice, HVAC control is difficult because outdoor temperature, solar gains, wind speed, and electricity prices all change over time. Building envelopes also differ from one site to another, so a controller must make good decisions under uncertainty and over long time horizons.

Conventional approaches range from simple rule-based or proportional--integral--derivative (PID) controllers to model predictive control (MPC). MPC can plan ahead by optimising over a building model, but it requires accurate calibration, forecasts, and nontrivial engineering effort, which limits scalability across buildings and climates. Deep reinforcement learning (DRL) offers a different path: a learned policy maps observations directly to control actions without solving a new optimisation problem at each step. This is attractive for HVAC because the task is sequential, continuous, and affected by delayed consequences. However, generalization remains a central issue, since a controller that performs well in one climate or simulator may lose performance when the deployment context changes.

This paper studies that question in a compact but informative setting. We use a single-zone thermal model with nonlinear conductance, solar gains, a coefficient-of-performance (COP)-aware energy term, and monthly electricity prices. We compare four continuous-control DRL agents: Proximal Policy Optimisation (PPO), Advantage Actor-Critic (A2C), Soft Actor-Critic (SAC), and Twin Delayed Deep Deterministic Policy Gradient (TD3). Training uses Syracuse weather only; evaluation uses Albany weather only, creating a strict geographic out-of-distribution test. We also add two training design choices aimed at improving generalization: domain randomisation over thermal parameters and round-robin seasonal sampling so that each agent sees winter, spring, summer, and fall evenly during training.

This paper makes three contributions. First, it provides a controlled comparison of four DRL agents on the same single-zone HVAC task with identical observations, actions, and reward structure. Second, it studies generalization with a clear train--test split across cities, not just across episodes from the same data source. Third, it shows that a lightweight simulator with physics-aware reward design, domain randomisation, and balanced seasonal exposure can produce policies that transfer well to unseen weather.

\begin{figure*}[!t]
\centering
\includegraphics[width=\textwidth]{pipeline.png}
\caption{End-to-end pipeline. EPA weather archives and New~York electricity
tariffs feed a nonlinear lumped-capacitance thermal simulator with domain
randomisation and stratified seasonal sampling. The simulator exposes a
Gymnasium environment with a 14-dimensional observation space, a continuous
HVAC power action, and a three-term reward balancing comfort, COP-adjusted
energy cost, and actuator smoothness. Four DRL agents (PPO, A2C, SAC, TD3)
are trained for 500\,k steps on Syracuse weather and evaluated
out-of-distribution on Albany over 100 deterministic episodes (25 per
season). Resulting episode logs and trained checkpoints are post-processed
into the comparative figures and tables of this paper.}
\label{fig:pipeline}
\end{figure*}

\section{Related Work}
\label{sec:related-work}
Deep reinforcement learning has attracted growing interest for HVAC control because it handles sequential decision-making, continuous action spaces, and delayed consequences without requiring explicit building models. Recent surveys confirm that DRL methods now match or exceed MPC on comfort--energy trade-offs, while highlighting open challenges in reward design, reproducibility, and sensitivity to changing conditions \cite{manjavacas2024drl_hvac,bekal2025continual_hvac,yu2025rl_urban_hvac}.

On the algorithm side, the actor-critic family is especially relevant for continuous control. A2C/A3C introduced asynchronous advantage estimation, which improved training stability through parallel data collection~\cite{mnih2016a2c}. PPO simplified trust-region policy updates with a clipped surrogate objective and became a widely used on-policy baseline~\cite{schulman2017ppo}. For off-policy control, TD3 reduced value overestimation with clipped double critics, delayed policy updates, and target policy smoothing~\cite{fujimoto2018td3}. SAC added maximum-entropy learning with a stochastic actor and strong sample efficiency, making it a popular choice for stable continuous control~\cite{haarnoja2018sac}.

Several recent studies compare these algorithms on HVAC tasks. The Sinergym benchmark~\cite{manjavacas2024drl_hvac} evaluated PPO, TD3, and SAC across different buildings and climates, finding that off-policy methods often achieve better comfort--energy trade-offs and arguing that many HVAC papers evaluate only one or two algorithms in a single narrow setting.

Generalization and transfer have emerged as a second major theme. A continual DRL study~\cite{bekal2025continual_hvac} showed that sample efficiency and transfer remain key bottlenecks; a model-based SAC framework with hypernetworks improved transfer across tasks and reduced catastrophic forgetting, suggesting that standard model-free training may not suffice when controllers must adapt across buildings. Along a complementary line, coupling DRL-based HVAC control with an urban climate model~\cite{yu2025rl_urban_hvac} demonstrated that background climate changes both the reward trade-off and the transferability of learned policies across cities.

Our work is closest in spirit to~\cite{manjavacas2024drl_hvac}, but focuses on a different gap. We use a simpler single-zone setting so that the effect of algorithm choice is easier to isolate, and we test transfer with a strict geographic split between training and evaluation cities. We also combine domain randomisation, balanced seasonal sampling, and a cost-aware reward that depends on both electricity price and instantaneous COP, bridging between high-level HVAC benchmarks and more controlled analysis of generalization.

\section{System Model}
\label{sec:system-model}
Fig.~\ref{fig:pipeline} gives an overview of the end-to-end system pipeline,
from raw data sources through the physics-based simulator and Gymnasium
environment to training and evaluation.

\subsection{Thermal Zone Dynamics}
\label{subsec:thermal-dynamics}
The indoor temperature evolves according to a first-law lumped-capacitance
model, discretised with a forward-Euler step of
$\Delta t = \SI{300}{\second}$:
%
\begin{equation}
  \label{eq:thermal-ode}
  T_{\text{in}}^{k+1} = T_{\text{in}}^{k}
    + \frac{\Delta t}{C}\Bigl[
        U_{\text{eff}}\!\bigl(T_{\text{out}}^{k} - T_{\text{in}}^{k}\bigr)
        + \alpha \cdot \text{GHI}^{k}
        - Q_{\text{hvac}}^{k}
      \Bigr]
\end{equation}
%
where the effective envelope conductance captures conduction, wind-driven
infiltration, and buoyancy-driven stack effects:
%
\begin{equation}
  \label{eq:ueff}
  U_{\text{eff}} = U_{\text{cond}}
    + k_w \, v_{\text{wind}}^{n_w}
    + k_s \sqrt{|T_{\text{in}} - T_{\text{out}}|}
\end{equation}
%
Table~\ref{tab:params} lists the physical parameters and their
domain-randomisation ranges.

\begin{table}[htb]
  \centering
  \caption{Thermal Simulator Parameters}
  \label{tab:params}
  \sisetup{per-mode=symbol}
  \begin{tabular}{llc}
    \toprule
    \textbf{Symbol} & \textbf{Description} & \textbf{Value / Range} \\
    \midrule
    $C$                  & Thermal capacitance       & $\mathcal{U}$(200\,k, 500\,k)\,\si{\joule\per\degreeCelsius} \\
    $U_{\text{cond}}$    & Conduction coefficient    & $\mathcal{U}$(30, 80)\,\si{\watt\per\degreeCelsius} \\
    $k_w$                & Wind infiltration coeff.  & \SI{4.45}{\watt\per\degreeCelsius} \\
    $n_w$                & Wind power-law exponent   & 0.65 \\
    $k_s$                & Stack effect coefficient  & \SI{2.6}{\watt\per\degreeCelsius} \\
    $\alpha$             & Solar gain coefficient    & \SI{0.5}{\metre\squared} \\
    $\eta_{\text{COP}}$  & Carnot efficiency factor  & 0.4 \\
    $Q_{\max}$           & Max HVAC power            & \SI{3000}{\watt} \\
    $T_{\text{set}}$     & Temperature setpoint      & \SI{22}{\degreeCelsius} \\
    $\Delta t$           & Control timestep          & \SI{300}{\second} \\
    \bottomrule
  \end{tabular}
\end{table}

\begin{table*}[!t]
  \centering
  \small
  \caption{14-Dimensional Observation Space (DR = domain randomisation;
  COP = Carnot coefficient of performance;
  $U_{\text{eff}}$ = effective envelope conductance).}
  \label{tab:obs}
  \sisetup{per-mode=symbol}
  \begin{tabular}{cllll}
    \toprule
    \textbf{\#} & \textbf{Feature} & \textbf{Range} & \textbf{Group} & \textbf{Description} \\
    \midrule
    0  & $T_{\text{in}} - T_{\text{set}}$  & $[-50,50]$\,\si{\degreeCelsius}  & Thermal   & Indoor temperature error \\
    1  & $(T_{\text{out}}-20)/20$           & $[-2,2]$                          & Thermal   & Outdoor temperature, centred and scaled \\
    2  & $\sin(2\pi h/24)$                  & $[-1,1]$                          & Time      & Hour sine encoding \\
    3  & $\cos(2\pi h/24)$                  & $[-1,1]$                          & Time      & Hour cosine encoding \\
    4  & $\sin(2\pi d/365)$                 & $[-1,1]$                          & Time      & Day sine encoding \\
    5  & $\cos(2\pi d/365)$                 & $[-1,1]$                          & Time      & Day cosine encoding \\
    6  & $p_t / p_{\max}$                   & $[0,1]$                           & Exogenous & Price normalised by annual maximum \\
    7  & GHI$/1200$                          & $[0,1]$                           & Exogenous & Solar irradiance normalised by 1200 \\
    8  & $v_{\text{wind}}/20$               & $[0,1]$                           & Exogenous & Wind speed normalised by 20 \\
    9  & $C / C_{\text{nom}}$               & $[0.57,1.43]$                     & DR        & Capacitance ratio \\
    10 & $U / U_{\text{nom}}$               & $[0.55,1.45]$                     & DR        & Conductance ratio \\
    11 & $a_{t-1}$                           & $[-1,1]$                          & Actuator  & Previous action \\
    12 & COP$/10$                            & $[0.08,1]$                        & Physics   & COP normalised by 10 \\
    13 & $U_{\text{eff}}/150$               & $[0,1]$                           & Physics   & Effective conductance normalised by 150 \\
    \bottomrule
  \end{tabular}
\end{table*}

\subsection{Weather Data and Train/Test Split}
\label{subsec:weather-data}
We use US Environmental Protection Agency (EPA) PRZM-format hourly archives from 1989 to 1990 for two
upstate New York stations:
\begin{itemize}
  \item \textbf{Syracuse} (station~14771), used for training only.
  \item \textbf{Albany} (station~14735), used for evaluation only as the \acOOD{} city.
\end{itemize}
Each archive provides about 17\,520 hourly records of temperature,
\acGHI{}, and wind speed. The two datasets have no overlapping records.
Albany data is never used during training, tuning, or model selection. The
split is geographic and enforced at the file-path level in the codebase, which
prevents data leakage by construction.

\subsection{Electricity Pricing}
\label{subsec:pricing}
A 12-element monthly electricity-price vector based on 2025 New York
residential rates is used to scale the energy term in the reward
function. The same price schedule is shared across training and
evaluation, since electricity pricing is a policy input rather than a
weather-dependent feature.

\section{Reinforcement Learning Formulation}
\label{sec:rl-formulation}

\subsection{Observation Space}
\label{subsec:state-space}
The agent observes a 14-dimensional continuous vector at each timestep,
summarised in Table~\ref{tab:obs}. Features are grouped into four
categories: thermal state, time encodings, pricing and weather, and
system parameters. Cyclic sine and cosine features encode hour of day and
day of year without discontinuities. Exposing the randomised ratios
$C/C_{\text{nom}}$ and $U/U_{\text{nom}}$ allows the policy
to adapt online to different building envelopes.

\subsection{Action Space}
\label{subsec:action-space}
The agent outputs a single continuous action $a_t \in [-1,\,1]$, mapped
to HVAC power by $Q_{\text{hvac}} = a_t \times Q_{\max}$
(Table~\ref{tab:action}). Positive values activate cooling. Negative
values activate heating. The continuous formulation lets the agent
modulate power smoothly instead of switching among discrete modes.

\begin{table}[!b]
  \centering
  \caption{Action Space}
  \label{tab:action}
  \begin{tabular}{ll}
    \toprule
    \textbf{Property} & \textbf{Value} \\
    \midrule
    Dimension       & 1 (\texttt{Box}) \\
    Range           & $a_t \in [-1,\,1]$ \\
    Mapping         & $Q_{\text{hvac}} = a_t \cdot Q_{\max}$ \\
    Max.\ power     & $Q_{\max} = \SI{3000}{\watt}$ \\
    \midrule
    $a_t > 0$       & Cooling \\
    $a_t < 0$       & Heating \\
    $a_t = 0$       & HVAC off \\
    \bottomrule
  \end{tabular}
\end{table}

\subsection{Reward Function}
\label{subsec:reward}
The agent receives a three-term penalty signal at each control step,
designed to balance thermal comfort, energy efficiency, and actuator
smoothness:
%
\begin{equation}
  \label{eq:reward}
  R_t = -\!\Bigl(
    \underbrace{w_c \,|T_{\text{in}} - T_{\text{set}}|}_{\text{comfort}}
    + \underbrace{w_e \,\frac{|a_t|}{\text{COP}_t} \cdot \frac{p_t}{p_{\max}}}_{\text{energy cost}}
    + \underbrace{w_s \,|a_t - a_{t-1}|}_{\text{slew penalty}}
  \Bigr)
\end{equation}
%
The comfort term penalises temperature error directly. The energy term
divides action magnitude by the instantaneous COP, so the same control
effort costs more when the system is less efficient, and then scales it by
time-of-use price. The slew penalty discourages abrupt changes in HVAC
power. The default weights are $w_c = 1.0$, $w_e = 0.1$, and
$w_s = 0.05$. This makes comfort the main objective while keeping the
controller cost-aware and smooth.

\subsection{Domain Randomisation}
\label{subsec:domain-rand}
At each episode reset the thermal capacitance $C$ and conduction
coefficient $U_{\text{cond}}$ are drawn uniformly from the ranges in
Table~\ref{tab:params}. The agent observes the ratios
$C/C_{\text{nom}}$ and $U/U_{\text{nom}}$, which allows it to adapt
to unknown building envelopes at test time.

\subsection{Stratified Seasonal Sampling}
\label{subsec:seasonal-sampling}
Naive random episode sampling from a two-year weather archive can
over-represent mild periods and under-sample the hardest heating and
cooling conditions. We therefore partition all valid 24-hour episode start
indices into four buckets: winter, spring, summer, and fall. At each reset,
the simulator cycles through these buckets in round-robin order and draws
a random start index from the active bucket.

The same protocol is used during evaluation. The 100 test episodes contain
exactly 25 episodes from each season. This keeps the reported metrics from
being dominated by whichever season contributes the most raw data.

\section{Algorithms}
\label{sec:algorithms}
We compare four continuous-action DRL algorithms implemented in
Stable-Baselines3~\cite{raffin2021sb3}: two on-policy methods,
PPO~\cite{schulman2017ppo} and A2C~\cite{mnih2016a2c}, and two
off-policy methods, SAC~\cite{haarnoja2018sac} and
TD3~\cite{fujimoto2018td3}.

\subsection{Proximal Policy Optimisation (PPO)}
\label{subsec:ppo}
PPO is an on-policy policy-gradient method that updates the actor with a
clipped surrogate objective. The clipping term limits how far the new
policy can move from the old policy in one update, which usually improves
stability. In our setting, PPO is a strong baseline because it is simple and
widely used, but its on-policy nature can make it less sample efficient than
off-policy methods.

\subsection{Advantage Actor-Critic (A2C)}
\label{subsec:a2c}
A2C belongs to the actor-critic family introduced through asynchronous
training by Mnih et al. It learns a policy and a value baseline together,
using advantage estimates to reduce gradient variance. Compared with PPO,
A2C is lighter and often faster to train, though it can be more sensitive to
hyperparameters and rollout settings.

\subsection{Soft Actor-Critic (SAC)}
\label{subsec:sac}
SAC is an off-policy actor-critic method based on the maximum-entropy
objective. It learns a stochastic policy that seeks high reward while also
keeping enough entropy for exploration. Replay-buffer reuse and entropy
regularisation make SAC an attractive choice for HVAC control, where data
efficiency and stable learning are both important.

\subsection{Twin Delayed Deep Deterministic Policy Gradient (TD3)}
\label{subsec:td3}
TD3 is an off-policy deterministic actor-critic method designed to reduce
value overestimation in continuous control. It uses two critics, delayed
policy updates, and target policy smoothing to stabilise learning. These
features are well suited to precise HVAC actuation, where small value errors
can lead to poor power commands.

\section{Experimental Setup}
\label{sec:experimental-setup}

\subsection{Training Configuration}
\label{subsec:training-config}
All agents are trained for 500\,000 environment steps on Syracuse
weather data. On-policy methods, PPO and A2C, use 8 parallel
\texttt{SubprocVecEnv} workers. Off-policy methods, SAC and TD3, use a
single \texttt{DummyVecEnv}. Observations are normalised via
\texttt{VecNormalize}. Reward normalisation is applied only for on-policy
agents. Table~\ref{tab:hyperparams} lists the key hyperparameters.

\begin{table}[htb]
  \centering
  \caption{Training Hyperparameters}
  \label{tab:hyperparams}
  \begin{tabular}{lcccc}
    \toprule
    \textbf{Parameter} & \textbf{PPO} & \textbf{A2C} & \textbf{SAC} & \textbf{TD3} \\
    \midrule
    Total steps         & \multicolumn{4}{c}{500\,000} \\
    Policy              & \multicolumn{4}{c}{MlpPolicy $[64,\,64]$} \\
    Seed                & \multicolumn{4}{c}{42} \\
    \midrule
    \# envs             & 8     & 8     & 1     & 1     \\
    Learning rate       & 3e-4  & 7e-4  & 3e-4  & 3e-4  \\
    $\gamma$            & 0.98  & 0.99  & 0.98  & 0.98  \\
    Batch size          & 64    & --    & 256   & 256   \\
    $n_{\text{steps}}$  & 2048  & 5     & --    & --    \\
    Buffer size         & --    & --    & 1\,M  & 1\,M  \\
    $\tau$              & --    & --    & 0.02  & 0.005 \\
    Entropy coeff.      & 0.01  & 0.01  & auto  & --    \\
    Action noise        & --    & --    & --    & $\mathcal{N}(0,0.1)$ \\
    Reward norm.        & yes   & yes   & no    & no    \\
    \bottomrule
  \end{tabular}
\end{table}

\subsection{Evaluation Protocol}
\label{subsec:eval-protocol}
Each trained agent is evaluated deterministically for 100 episodes on
the unseen Albany weather data, with 25 episodes from each season. We
report the median of five metrics: mean absolute error (MAE), root-mean-square error (RMSE), time within
$\pm\SI{0.5}{\degreeCelsius}$ of setpoint, daily energy consumption, and
peak HVAC power.

\begin{figure}[!b]
  \centering
  \includegraphics[width=\columnwidth]{training_convergence.png}
  \caption{Mean evaluation reward during training (500k steps, 8 parallel
  environments for PPO/A2C, 1 for SAC/TD3). Shaded regions show the
  interquartile range across evaluation episodes.}
  \label{fig:convergence}
\end{figure}

\section{Results}
\label{sec:results}

\subsection{Learning Curves}
\label{subsec:learning-curves}
Fig.~\ref{fig:convergence} shows the mean evaluation reward during training.
TD3 and SAC converge smoothly and plateau by 400k steps. A2C is competitive
despite early instability. PPO reaches a lower reward plateau, consistent
with its weaker test-time performance.

\subsection{Generalization Performance}
\label{subsec:generalisation}
Table~\ref{tab:results} summarises the overall Albany test-set
performance. TD3 achieves the lowest MAE (\SI{0.12}{\degreeCelsius})
and the lowest energy use (\SI{16.9}{\kilo\watt\hour\per day}) while
maintaining near-perfect comfort with $99.7\%$ of time within
$\pm\SI{0.5}{\degreeCelsius}$ of setpoint. A2C matches TD3 on comfort
but consumes slightly more energy. PPO is the weakest method on all main
metrics and reaches the comfort band only $64.9\%$ of the time.

\begin{table}[htb]
  \centering
  \caption{Median Albany Performance (100 OOD Episodes)}
  \label{tab:results}
  \begin{tabular}{lcccc}
    \toprule
    \textbf{Metric}
      & \textbf{A2C} & \textbf{PPO} & \textbf{SAC} & \textbf{TD3} \\
    \midrule
    MAE (\si{\degreeCelsius})
      & 0.16 & 0.45 & 0.17 & \textbf{0.12} \\
    RMSE (\si{\degreeCelsius})
      & 0.19 & 0.50 & 0.22 & \textbf{0.15} \\
    Within $\pm$0.5\,\si{\degreeCelsius} (\%)
      & \textbf{99.7} & 64.9 & 99.0 & \textbf{99.7} \\
    Energy (\si{\kilo\watt\hour/day})
      & 18.5 & 21.7 & 21.0 & \textbf{16.9} \\
    Peak Power (\si{\watt})
      & 2657 & 3000 & \textbf{2415} & 2397 \\
    \bottomrule
  \end{tabular}
\end{table}

Fig.~\ref{fig:albany-heatmap} breaks performance down by season.
All algorithms struggle most in winter, where the larger indoor to outdoor
temperature difference increases both tracking difficulty and energy use.
TD3 and A2C keep at least $99.7\%$ comfort across all four seasons,
while PPO drops sharply in spring.

\begin{figure}[htb]
  \centering
  \includegraphics[width=\columnwidth]{albany_ood_summary.png}
  \caption{Median MAE (\si{\degreeCelsius}) on Albany (out-of-distribution)
  over 25 episodes per season. Lower is better.}
  \label{fig:albany-heatmap}
\end{figure}

Fig.~\ref{fig:gap} shows how performance changes when moving from the
Syracuse training distribution to the unseen Albany test distribution.
Comfort degradation is very small for A2C, SAC, and TD3, while PPO shows
a much larger drop. All agents also use less energy on Albany, which is
consistent with a milder test climate.

\begin{figure}[htb]
  \centering
  \includegraphics[width=\columnwidth]{generalization_gap.png}
  \caption{Percentage change in median performance when transferring
  from Syracuse (training) to Albany (test). Positive values indicate
  worse transfer performance on Albany; comfort is reported as loss in
  percentage points.}
  \label{fig:gap}
\end{figure}

\subsection{Seasonal Analysis}
\label{subsec:seasonal}
Fig.~\ref{fig:traces} shows representative winter and summer episodes
for TD3, the best-performing algorithm overall. These seasons were chosen
because they reflect the most demanding heating and cooling conditions. In
winter, the controller applies sustained heating of about
\SI{2000}{\watt} to counter a $-10\,^{\circ}\text{C}$ outdoor temperature
while keeping the indoor temperature inside the comfort band. In summer,
less HVAC effort is needed and the policy changes power smoothly.

\begin{figure*}[!t]
  \centering
  \includegraphics[width=\textwidth]{best_policy_traces.png}
  \caption{Representative winter and summer Albany episodes for TD3.
  Left: indoor and outdoor temperature traces with the
  $\pm\SI{0.5}{\degreeCelsius}$ comfort band in green.
  Right: HVAC power profile for each season.}
  \label{fig:traces}
\end{figure*}

\subsection{Discussion}
\label{subsec:discussion}
Three patterns stand out. First, TD3 is the best overall method in this
setting. Its deterministic policy and conservative critic updates appear to
fit a low-dimensional continuous control problem where precise power levels
matter. The gains are visible in both comfort and energy use.

Second, A2C is a stronger baseline than expected. Although it is simpler
than SAC and TD3, it reaches the same comfort rate as TD3 and stays close
on tracking error. This suggests that the environment is structured enough
for a lightweight on-policy method to learn a good policy when training is
stable and season coverage is balanced.

Third, the transfer results suggest that the training protocol matters as
much as the algorithm. Domain randomisation exposes each agent to a range
of thermal masses and conductances. Stratified seasonal sampling prevents
the policy from overfitting to mild conditions. We cannot isolate the exact
contribution of each design choice from the current experiments, but the
small comfort gap for A2C, SAC, and TD3 on the unseen city is consistent
with both choices improving transfer performance.

PPO is the main outlier. Its lower comfort rate and higher energy use point
to weaker policy updates for this task under the current budget and
hyperparameters. This does not mean PPO is unsuitable for HVAC control in
general, but it was clearly less effective here than the other three agents.

\section{Conclusion}
\label{sec:conclusion}
This paper presented a controlled comparison of four DRL agents---PPO, A2C,
SAC, and TD3---for continuous single-zone HVAC control under a common
nonlinear thermal simulator, reward function, and evaluation protocol.
All agents were trained exclusively on Syracuse weather and evaluated on
100 deterministic episodes of unseen Albany weather (25 per season),
providing a strict geographic out-of-distribution test.

TD3 achieved the strongest overall performance on the Albany test set,
recording a median MAE of \SI{0.12}{\degreeCelsius}, an RMSE of
\SI{0.15}{\degreeCelsius}, $99.7\%$ of time within
$\pm\SI{0.5}{\degreeCelsius}$ of the \SI{22}{\degreeCelsius} setpoint,
and daily energy use of only \SI{16.9}{\kilo\watt\hour\per day} with a
peak draw below \SI{2400}{\watt}. A2C matched TD3 on comfort ($99.7\%$
in-band) with a competitive MAE of \SI{0.16}{\degreeCelsius}, though at a
higher daily cost of \SI{18.5}{\kilo\watt\hour\per day}. SAC maintained
$99.0\%$ comfort and the lowest peak power (\SI{2415}{\watt}), but its
energy use of \SI{21.0}{\kilo\watt\hour\per day} was higher than TD3.
PPO was the weakest agent, with an MAE of \SI{0.45}{\degreeCelsius}, only
$64.9\%$ comfort, the highest energy consumption at
\SI{21.7}{\kilo\watt\hour\per day}, and actuator saturation at the
\SI{3000}{\watt} system limit.

The generalization gap analysis confirmed that A2C, SAC, and TD3 all
transferred well from Syracuse to Albany, with comfort losses below one
percentage point. All agents used less energy on Albany than on Syracuse,
consistent with Albany's milder climate. PPO was the outlier, showing a
substantially larger comfort drop on the unseen city.

Two training-design choices contributed to these results. Domain
randomisation over thermal capacitance~$C$ and conductance~$U$ forced each
agent to handle a range of building envelopes, while round-robin seasonal
sampling ensured balanced exposure to winter heating and summer cooling
extremes. Although the current experiments do not isolate the individual
contribution of each mechanism, the small generalization gap for three of
four agents supports their combined effectiveness.

Several directions remain open. Extending the simulator to multi-zone
buildings with inter-zone heat flow would test scalability. Adding
occupancy schedules, demand-response events, and time-varying tariffs would
increase realism. Formal ablation studies on domain randomisation and
seasonal sampling are needed to quantify their separate effects. Finally,
increasing the training budget beyond 500\,k steps or tuning
PPO-specific hyperparameters may close the gap between PPO and the
stronger agents.

\bibliographystyle{IEEEtran}
\bibliography{references}

\end{document}
